\clearpage
\section*{Anexo 4-V — Protocolos de aceptación}
\phantomsection\label{fuente-25-anexo-4-1-v-aceptacion-y-dependencias}
\phantomsection
\label{anx:V}
\addcontentsline{toc}{section}{Anexo 4-V — Protocolos de aceptación}
Las pruebas siguientes se ejecutan en Preproducción y antes de la ola aplicable. Cada acta identifica carga, inyección de falla, resultados observados y responsables. El Anexo 4-M desarrolla AL-DTE-01, AL-DR-01, AL-OFF-01 y AL-CLI-01. La Tabla~\ref{tab:condiciones-aceptacion} resume sus protocolos.

\begin{tablalafrox}{Protocolos de aceptación de la arquitectura lógica}{tab:condiciones-aceptacion}{F{0.17} F{0.29} F{0.33} Y}{\cab{ID} & \cab{Ejecución} & \cab{Criterio de aceptación} & \cab{Evidencia}}
\mbox{AL-DTE-01} & Ensayar 96 guías por las rutas de Talca y de los demás sitios, con fallas y cambio de carga. & Cada camión sale con guía válida del ERP, sin los tres caminos, solo lleva la carga amparada. Una guía anulada no se reutiliza. & Origen, ruta, folio, carga, acuse y acta. \\
\mbox{AL-POD-01} & Ensayar entrega total, parcial y receptor sustituto sin señal. & El POD conserva autor, objeto y acuse aplicable. & Firma, hash y trazas. \\
\mbox{AL-SLA-01} & Simular ventanas y cuotas de ERP, SII y contrapartes. & Reintentos y escalamiento respetan cada contrato. & Contratos, logs y tiempos. \\
\mbox{AL-OFF-01} & Aislar cada sitio 24 horas con dos relevos y despachar con guía preemitida. & Operación local y conciliación sin duplicados, solo sale carga amparada. & Bitácora, guías, colas y acta. \\
\mbox{AL-ID-01} & Ensayar reloj, expiración y baja conectada/aislada; ruta de 14,5 h. & Cero accesos indebidos; reparto acotado a 16 h; cola conservada. & Manifiesto, política, tiempos; Anexo 4-M. \\
\mbox{AL-CLI-01} & Armar pedidos sin señal en el Portal de Clientes instalado y reconectar. & Ningún pedido se pierde ni se duplica, M3 confirma solo al recibirlo y obtener la retención de M2 y el control de crédito de M7. & UUID, respuesta de M3 y acta. \\
\mbox{AL-DR-01} & Cortar caminos y recuperar Talca en nube. & RTO $\leq$ 4 horas y RPO $\leq$ 15 minutos. & Marcas de agua y acta. \\
\mbox{AL-PERF-01} & Aplicar carga de 3$\times$ con mamparos por perfil. & Latencias y colas cumplen el Anexo 4-T. & Métricas y reporte. \\
\mbox{AL-STOCK-01} & Pedir la última unidad desde preventa, portal y EDI, con cortes y acuses perdidos. & Una sola reserva confirmada con retención durable, nunca stock negativo. & Estado central y local, p95 y acta. \\
\mbox{AL-ACT-01} & Probar cada actor del Anexo 4-N con una operación permitida y una denegada. & Autorización por acción y recurso, con auditoría personal. & Política, etapa, registros y acta. \\
\mbox{AL-MOD-01} & Revisar contratos y ADR, reiniciar consumidor y repetir UUID. & Un efecto durable, dueño único, cero confirmaciones ficticias. & Contrato, oráculo y bitácora. \\
\mbox{AL-FRIO-01} & Inyectar excursión y liberación por roles distintos, con corte y reenvío. & Solo Calidad decide; la retención persiste ante permiso o evidencia inválidos. & Regla, versión, actor y estado. \\
\mbox{AL-UI-01} & Revisar accesibilidad y exportación de auditoría por perfil. & WCAG A/AA sin hallazgos abiertos; exportación completa y autorizada. & Informe, filtros y conjunto esperado. \\
\end{tablalafrox}
{\fontsize{9pt}{11pt}\selectfont\raggedright Fuente: caso, RT-03.10/13 (autonomía y sincronización) y cap. 18 (aceptación); Transversales, RT-02.06, RT-11.01 y RT-12.05/07 (idempotencia y autorización). Protocolos propuestos LafroX de anexos 4-M y 4-V.\par}

AL-DTE-01 comprueba ambas rutas de las 96 guías, la preemisión nocturna, la invalidación por cambio de carga y la nueva emisión del ERP ante el SII por fibra, LTE o Starlink en espera caliente. AL-OFF-01 comprueba por separado el despacho con guía válida preemitida. AL-DR-01 mide la extracción continua y la recuperación del WMS de Talca en ECS Fargate. El esquema de replicación y continuidad se desarrolla en 4.3.2. Los resultados de todas las pruebas se vinculan a los requisitos y al registro ADR del Anexo 4-O.

\subsection*{AL-STOCK-01. Confirmación central con custodia local}
En Preproducción se pide la misma última unidad al mismo tiempo desde preventa, portal y EDI. Se corta el enlace antes y después del registro local, se pierde el acuse, se repite el UUID, se cancela durante un aislamiento y se presenta una época obsoleta. Solo se acepta una reserva confirmada con retención durable de igual identidad, época y cantidad. El disponible nunca queda negativo y un timeout nunca confirma. La preparación consume la retención y una cancelación tardía no libera lo consumido. La prueba mide el p95 de confirmación con la latencia entre la nube y el sitio, la edad de la cola y la recuperación. También mide la proporción de líneas que requieren más de una retención, para contrastar el volumen del Anexo 4-I. Los resultados se trazan a RNG-01, RF-03.03, RT-02.06 y RT-03.12.

\subsection*{AL-ACT-01. Autorización de los quince actores}
\enlargethispage{2\baselineskip}
Cada fila del Anexo 4-N prueba una operación permitida y otra denegada. Se incluyen el cruce entre empresas o rutas, la credencial de un representante usada como conductor, un rol de TI que intenta una aprobación de negocio, la aprobación del descuadre propio, la baja y el relevo sin conexión. Se comprueba que cada portal se habilita solo en su etapa, que el cliente tradicional compra asistido sin cuenta y que el catálogo público no muestra precios. El resultado esperado es una autorización por acción y recurso, con auditoría personal y sin credenciales compartidas. La evidencia se registra por actor, requisito, versión de la política y etapa, y se traza a RT-12.05, RT-12.06 y RT-12.12.


\subsection*{AL-MOD-01. Fronteras, reintentos y degradación lógica}
Se eligen los recorridos pedido--reserva--preparación y reparto--rendición. La revisión de contratos comprueba dueño de cada entidad, interfaz pública, versión y ADR aplicable; deptrac debe rechazar una dependencia no autorizada entre contextos. Se repite una solicitud con el mismo UUID antes y después del commit y se pierde su acuse: debe recuperarse el resultado original sin duplicar escritura ni evento. Se reinicia el consumidor entre recepción y confirmación y se reenvían eventos fuera de orden; el estado final debe coincidir con el oráculo y toda excepción queda auditada. La caída del ERP no permite simular una guía emitida; el reintento conserva la correlación. Para stock se aplica además AL-STOCK-01. La aceptación exige cero efectos duplicados, cero escrituras por un módulo no dueño y ningún estado confirmado sin su condición durable.

\subsection*{AL-FRIO-01. Bloqueo y liberación autorizada}
Se carga una regla térmica de prueba versionada y un SSCC de cada familia ensayada; se inyectan excursión crítica, datos incompletos y nueva excursión después de una liberación. Con y sin WAN, M2/M5 mantienen la retención y M6 no habilita entrega del SSCC afectado. Un conductor y un administrador TI intentan liberar: se deniega. Calidad con evidencia suficiente y permiso vigente decide; se comprueba que M9 guarda una sola disposición y que M2/M5 aplican su misma versión, incluso tras reinicio o reenvío. Permiso vencido o expediente incompleto no liberan. Se conserva estado anterior/posterior, regla, actor y trazas. El protocolo comprueba autoridad lógica; no acredita calibración de sensores ni ensayos vehiculares.

\subsection*{AL-UI-01. Acceso humano y auditoría autorizada}
En pedido, picking, entrega y Calidad se comprueba WCAG 2.2 AA de forma automática y manual: teclado, foco, contraste, nombres y errores comprensibles. Ningún hallazgo A/AA confirmado queda abierto al aceptar. En auditoría, combinar filtros de persona, período, entidad y operación; CSV/JSON debe incluir todo el resultado autorizado, no solo la página visible. Una cuenta de otra empresa no consulta ni exporta esos datos. El expediente conserva perfil, versión, conjunto esperado, resultado y correcciones. Es un protocolo propuesto, no un ensayo ejecutado.

\label{fuente-25-anexo-4-1-v-aceptacion-y-dependencias-fin}
